\section{The TW128 Construction}
\label{sec:tw128}

\subsection{Tree Structure}
\label{sec:design}

\TW{} is a nonce-based AEAD built as a two-level tree of duplex objects over
\Keccakp{}. The \emph{root} duplex is initialized with $(K, U)$, absorbs the
associated data and the first ciphertext chunk, and emits the final
$\tauroot$-bit root tag $T$. The remaining plaintext is partitioned into
\emph{leaf} chunks $M_1,\dots,M_n$; leaf $j$ is encrypted by an independent
duplex initialized with $(K, U, j)$ and returns a $\tauleaf$-bit leaf tag $T_j$.
For multi-chunk messages, the root aggregates
$T_1 \concat \cdots \concat T_n \concat \langle n \rangle_8$ before emitting
$T$. Thus the authentication tag is also the digest point for the full
$(K, U, A, M)$ context. The construction is depicted in \Cref{fig:tw128-tree};
the exact edge-case conventions are specified in \Cref{sec:enc}.

The tree fixes the transcript but not the implementation's parallelism. Leaves
operate on disjoint states after $(K, U)$ is fixed, so scalar, SIMD, and
multicore implementations can evaluate different numbers of leaves at once
while producing identical ciphertexts and tags. The root state and each leaf
state are bounded $1600$-bit objects, and leaf tags can be aggregated as a
stream, giving constant working memory at any fixed degree of parallelism. This
is the AEAD analog of KangarooTwelve's final-node hashing of per-chunk chaining
values, with leaf tags as the chaining values and the root tag as the final
digest.

\begin{figure}[t]
  \centering
  \begin{tikzpicture}[
    font=\small,
    >=Stealth,
    block/.style={draw, rounded corners=1pt, minimum height=10mm,
                  align=center},
    init/.style={block, minimum width=15mm, fill=black!10,
                 font=\footnotesize},
    msg/.style={block, minimum width=11mm, fill=blue!12},
    cv/.style={block, minimum width=9mm, fill=green!18,
               font=\footnotesize},
    enc/.style={block, minimum width=11mm, fill=black!10,
                font=\scriptsize},
    leaf/.style={block, minimum height=22mm, minimum width=11mm,
                 fill=blue!12},
    arr/.style={->, thin},
    flab/.style={font=\scriptsize, inner sep=1pt},
  ]
    \node[init] (init) {$K, U$};
    \node[msg, right=1mm of init] (ad) {$A$};
    \node[msg, right=1mm of ad] (c0) {$C_0$};
    \node[cv, right=1mm of c0] (t1) {$T_1$};
    \node[cv, right=1mm of t1] (t2) {$T_2$};
    \node[right=1mm of t2, inner sep=1pt] (dots) {$\cdots$};
    \node[cv, right=1mm of dots] (tn) {$T_n$};
    \node[enc, right=1mm of tn] (lenc) {$\langle n \rangle_8$};

    \draw[arr] (lenc.east) -- ++(7mm,0) node[right] {$T$};

    \node[leaf, above=12mm of t1] (leaf1) {$C_1$};
    \node[leaf, above=12mm of t2] (leaf2) {$C_2$};
    \node[leaf, above=12mm of tn] (leafn) {$C_n$};

    \node[flab, above=1mm of leaf1.north] {$K, U, 1$};
    \node[flab, above=1mm of leaf2.north] {$K, U, 2$};
    \node[flab, above=1mm of leafn.north] {$K, U, n$};

    \draw[arr] (leaf1.south) -- (t1.north);
    \draw[arr] (leaf2.south) -- (t2.north);
    \draw[arr] (leafn.south) -- (tn.north);
  \end{tikzpicture}
  \caption{Schematic of \TW{} for $n \geq 1$ leaves. The root duplex
    (bottom row) absorbs the associated data, root ciphertext chunk, leaf tags,
    and leaf count before emitting the root tag $T$. Each independent leaf
    duplex (top row) encrypts one leaf chunk and emits a leaf tag.}
  \label{fig:tw128-tree}
\end{figure}

\subsection{Parameters and Domain}
\label{sec:params}

\Cref{tab:params} lists the user-visible and derived scheme parameters.
The implementation uses the same byte length for the root tag and each
leaf tag, but the proofs of \Cref{sec:aead-security,sec:root-tag-hash-security} keep
$\tauroot$ and $\tauleaf$ as separate symbols.

\begin{table}[ht]
  \centering
  \small
  \begin{tabular}{l l l}
    \toprule
    Parameter & Value & Description \\
    \midrule
    $k$                  & $256$          & Key length (bits) \\
    $\nu$                & $256$          & Nonce length (bits) \\
    $\tauroot$ & $256$          & Root tag length (bits) \\
    $\tauleaf$ & $256$          & Leaf tag length (bits) \\
    $\beta$              & $8183$         & Plaintext chunk size (bytes) \\
    $M_{\max}$           & $8183 \cdot 2^{64}$ bytes & Maximum plaintext length \\
    $A_{\max}$           & $8183 \cdot 2^{64}$ bytes & Maximum associated data length \\
    \midrule
    $b$           & $1600$  & Permutation state width (bits) \\
    $n_r$         & $12$    & Permutation rounds \\
    $r$           & $1344$  & Sponge rate (bits) \\
    $c$           & $256$   & Sponge capacity (bits) \\
    $\rho$        & $167$   & Maximum byte-aligned absorption per duplex call (bytes) \\
    \bottomrule
  \end{tabular}
  \caption{\TW{} parameters.}
  \label{tab:params}
\end{table}

The bound $M_{\max} = 8183 \cdot 2^{64}$ bytes is the largest plaintext whose leaf
index fits in $\langle \cdot \rangle_8$: writing $\leaves(M) =
\max(0, \lceil \|M\| / \beta \rceil - 1)$, every
supported plaintext satisfies $0 \leq \leaves(M) \leq 2^{64} - 1$, so
leaf indices lie in $\{1, \dots, 2^{64} - 1\}$. The associated data
bound $A_{\max}$ is a fixed scheme parameter, taken equal to
$M_{\max}$. Inputs exceeding these bounds are outside the
scheme domain.

The derived per-call absorption bound $\rho$ follows the \cite{BDPV11}
analysis: with $r = 1344$ and $\padtenone$ padding, $\rho_{\max}(\padtenone, r) =
r - 2 = 1342$ bits. The byte-aligned data field shares this allotment
with the four-bit phase suffix of \Cref{sec:domain-sep}, so $\rho$ is
the largest integer satisfying $8\rho + 4 \leq 1342$, giving $\rho =
167$ data bytes (and $8\rho + 4 = 1340 \leq 1342$).

\paragraph{Parameter rationale.}
The capacity, key length, and tag lengths are all set to $256$ bits so the
flat sponge loss, key-guessing terms, tag-guessing terms, and root/leaf
collision terms meet a $128$-bit target at the concrete work cap of
\Cref{sec:concrete-claim}. With $b=1600$, this fixes the rate at
$r=1344$. The 12-round \Keccakp{} permutation matches KangarooTwelve
\cite{BDPVVV16}; \Cref{sec:permutation-cryptanalysis} discusses the
cryptanalytic record for this round count. The chunk size is
$\beta = 49\rho = 8183$ bytes, the largest multiple of the byte-aligned
per-call absorption bound not exceeding $2^{13}$ bytes; this fully utilizes
message duplex calls inside full leaves while keeping the leaf working set
small. The $256$-bit nonce is not required by the nonce-respecting bounds, but
it is essentially free in the root initialization call and supports both random
nonces and narrower counter nonces by zero-padding.

\subsection{Domain Separation}
\label{sec:domain-sep}

Distinct duplex instances and phases are separated at the transcript level by
initialization prefixes and phase suffixes. The initialization strings carry a
6-byte ASCII prefix and, for leaves, an eight-byte little-endian chunk index:
\begin{align*}
  \prfxroot &= \texttt{"TW128R"} \in \bits^{48}, \\
  \prfxleaf &= \texttt{"TW128L"} \in \bits^{48}.
\end{align*}
The root's first absorbed string is $\prfxroot \concat K \concat U$, and leaf
$j$'s first absorbed string is
$\prfxleaf \concat K \concat U \concat \langle j \rangle_8$. Every
byte-aligned absorption also carries one of the four-bit phase suffixes in
\Cref{tab:suffixes}, written least significant bit first: absorbing a byte
block $\sigma$ under suffix \sigph{} feeds
$\sigma \concat \sigph$ to the duplex interface of \Cref{sec:sponge} before
$\padtenone$ padding.

\begin{table}[ht]
  \centering
  \small
  \begin{tabular}{l l l}
    \toprule
    Name & Hex & Use \\
    \midrule
    \siginit & \texttt{0xC} & End of initialization (the sole init call) \\
    \sigadmore   & \texttt{0x9} & Non-final associated data block \\
    \sigadlast   & \texttt{0xD} & Final associated data block \\
    \sigmsgmore  & \texttt{0xA} & Non-final message block \\
    \sigmsglast  & \texttt{0xE} & Final message block \\
    \sigaggmore  & \texttt{0xB} & Non-final aggregation block \\
    \sigagglast  & \texttt{0xF} & Final aggregation block (emits root tag) \\
    \bottomrule
  \end{tabular}
  \caption{Phase suffixes (4 bits each, appended LSB-first).}
  \label{tab:suffixes}
\end{table}

\Cref{tab:transcript-schedule} summarizes the root and leaf transcript
schedule. The algorithms of \Cref{sec:enc,sec:dec} give the normative
byte-level specification; the table records which phase is present and what
output the final call of that phase supplies.

\begin{table}[ht]
  \centering
  \scriptsize
  \setlength{\tabcolsep}{1pt}
  \begin{tabular}{@{}p{0.10\linewidth}p{0.11\linewidth}p{0.13\linewidth}p{0.26\linewidth}p{0.24\linewidth}p{0.09\linewidth}@{}}
    \toprule
    Instance & Phase & Suffixes & Absorbed string & Output supplied & When \\
    \midrule
    Root
      & Init
      & \siginit
      & $\prfxroot \concat K \concat U$
      & $0$ if $A \neq \varepsilon$; otherwise first $M_0$ keystream
      & Always \\
    Root
      & AD
      & \sigadmore, \sigadlast
      & $\rho$-byte block partition of $A$
      & Final call emits first $M_0$ keystream
      & $A \neq \varepsilon$ \\
    Root
      & Message
      & \sigmsgmore, \sigmsglast
      & $\rho$-byte block partition of $C_0$
      & Non-final calls emit next keystream; final call emits $T$ if $n = 0$,
        and $0$ otherwise
      & Always \\
    Leaf $j$
      & Init
      & \siginit
      & $\prfxleaf \concat K \concat U \concat \langle j \rangle_8$
      & First $M_j$ keystream
      & $1 \leq j \leq n$ \\
    Leaf $j$
      & Message
      & \sigmsgmore, \sigmsglast
      & $\rho$-byte block partition of $C_j$
      & Non-final calls emit next keystream; final call emits $T_j$
      & $1 \leq j \leq n$ \\
    Root
      & Aggregation
      & \sigaggmore, \sigagglast
      & $T_1 \concat \cdots \concat T_n \concat \langle n \rangle_8$
      & Final call emits $T$
      & $n \geq 1$ \\
    \bottomrule
  \end{tabular}
  \caption{\TW{} transcript schedule. ``First $M_i$ keystream'' means
    $8\min(\|M_i\|,\rho)$ output bits for the first duplex block of that
    message chunk; subsequent message-phase outputs have the length of the next
    duplex block, as in \Cref{alg:helpers}.}
  \label{tab:transcript-schedule}
\end{table}

The combined effect of the initialization prefixes and the phase
suffixes is that the flattened sponge input of any \TW{} duplex call
uniquely determines the duplex instance, the leaf index when applicable,
the phase, and the duplex-call position within that phase.
\Cref{def:wf-transcripts} (Well-formed \TW{} transcript prefixes),
\Cref{lem:transcript-inj} (Transcript Injectivity), and its structural
corollaries formalize this parsing fact for the proofs of
\Cref{sec:aead-security,sec:root-tag-hash-security} and
\Cref{cor:committing-combined}.

\subsection{Encryption}
\label{sec:enc}

\Cref{alg:enc} specifies $\Enc_K(U, A, M)$ using the helper procedures of
\Cref{alg:helpers}. The helpers sit below the tree layer: the encryption and
decryption algorithms split bodies into $\beta$-byte chunks, while the helpers
split a chosen byte string into $\rho$-byte duplex blocks. \Call{StreamEnc}{}
XORs each plaintext block with the current keystream, absorbs the resulting
ciphertext block under the message suffix, and uses the duplex output as the
next keystream. This is the duplex-level presentation of the overwrite
mechanism of \cite[\S6.2 and Algorithm~5]{BDPV11}. The decryption helper
\Call{StreamDec}{} uses the same schedule with ciphertext blocks as input:
it recovers $m_i = c_i \oplus Z$, absorbs the same $c_i$, and returns the
plaintext string and final output value.

\paragraph{Edge-case convention.}
The helper partition always has at least one block: an empty input is the single
empty final block carrying the final phase suffix. The associated-data phase is
present exactly when $A \neq \varepsilon$; when $A = \varepsilon$, the root
\siginit call emits the first root-chunk keystream. The first keystream request
for a message string $X$ is $8\min(\|X\|, \rho)$ bits. If $n = 0$, the root
message phase requests $\tauroot$ bits on its final \sigmsglast call and no
aggregation phase occurs. If $n \geq 1$, that root message call requests zero
bits, each leaf final message call requests $\tauleaf$ bits, and the final
\sigagglast aggregation call requests $\tauroot$ bits. These conventions apply
identically to decryption.

\begin{algorithm}[t]
  \caption{TW128 helper procedures.}
  \label{alg:helpers}
  \begin{algorithmic}[1]
    \Procedure{Absorb}{$\mathcal{D}$, $s$, $\sigma^-$, $\sigma^+$, $\ell$}
      \State Partition $s$ into duplex blocks $s_1, \dots, s_t$ of exactly
        $\rho$ bytes each, except that the last block may be shorter.
      \State Require $t \geq 1$; the empty case yields the single empty
        duplex block $s_1 = \varepsilon$.
      \For{$i \gets 1$ \textbf{to} $t - 1$}
        \State $\mathcal{D}.\duplexing(s_i \concat \sigma^-, 0)$
      \EndFor
      \State \Return $\mathcal{D}.\duplexing(s_t \concat \sigma^+, \ell)$
    \EndProcedure
    \Statex
    \Procedure{StreamEnc}{$\mathcal{D}$, $m$, $Z$, $\sigma^-$, $\sigma^+$, $\ell$}
      \State Partition $m$ into duplex blocks $m_1, \dots, m_t$ of exactly
        $\rho$ bytes each, except that the last block may be shorter.
      \State Require $t \geq 1$; if $m = \varepsilon$, set $t = 1$ and
        $m_1 = \varepsilon$.
      \For{$i \gets 1$ \textbf{to} $t - 1$}
        \State $c_i \gets m_i \oplus Z$
        \State $Z \gets \mathcal{D}.\duplexing(c_i \concat \sigma^-, 8\|m_{i+1}\|)$
      \EndFor
      \State $c_t \gets m_t \oplus Z$
      \State $Z' \gets \mathcal{D}.\duplexing(c_t \concat \sigma^+, \ell)$
      \State \Return $(c_1 \concat \cdots \concat c_t, Z')$
    \EndProcedure
  \end{algorithmic}
\end{algorithm}

\begin{algorithm}[t]
  \caption{$\Enc_K(U, A, M)$.}
  \label{alg:enc}
  \begin{algorithmic}[1]
    \Require $K \in \bits^k$, $U \in \bits^{\nu}$, $A$ with $\|A\| \leq A_{\max}$, $M$ with $\|M\| \leq M_{\max}$.
    \Ensure $C \concat T$ with $\|C\| = \|M\|$ and $T \in \bits^{\tauroot}$.
    \State $M_0 \gets$ first $\min(\|M\|, \beta)$ bytes of $M$ (possibly $\varepsilon$).
    \State Partition the remaining suffix into chunks $M_1, \dots, M_n$ of exactly $\beta$ bytes each, except that the last chunk may be shorter but nonempty, so $n = \leaves(M)$.
    \State $\mathcal{D}_{\text{root}} \gets$ new $\Duplex$
    \If{$A = \varepsilon$}
      \State $Z \gets \mathcal{D}_{\text{root}}.\duplexing(\prfxroot \concat K \concat U \concat \siginit, 8\min(\|M_0\|, \rho))$
    \Else
      \State $\mathcal{D}_{\text{root}}.\duplexing(\prfxroot \concat K \concat U \concat \siginit, 0)$
      \State $Z \gets \Call{Absorb}{\mathcal{D}_{\text{root}}, A, \sigadmore, \sigadlast, 8\min(\|M_0\|, \rho)}$
    \EndIf
    \If{$n = 0$}
      \State $(C_0,\, T) \gets \Call{StreamEnc}{\mathcal{D}_{\text{root}}, M_0, Z, \sigmsgmore, \sigmsglast, \tauroot}$
      \State \Return $C_0 \concat T$
    \EndIf
    \State $(C_0,\, \_) \gets \Call{StreamEnc}{\mathcal{D}_{\text{root}}, M_0, Z, \sigmsgmore, \sigmsglast, 0}$
    \For{$j \gets 1$ \textbf{to} $n$ \textbf{in parallel}}
      \State $\mathcal{D}_j \gets$ new $\Duplex$
      \State $Z_j \gets \mathcal{D}_j.\duplexing(\prfxleaf \concat K \concat U \concat \langle j \rangle_8 \concat \siginit, 8\min(\|M_j\|, \rho))$
      \State $(C_j, T_j) \gets \Call{StreamEnc}{\mathcal{D}_j, M_j, Z_j, \sigmsgmore, \sigmsglast, \tauleaf}$
    \EndFor
    \State $G \gets T_1 \concat \cdots \concat T_n$.
    \State $T \gets \Call{Absorb}{\mathcal{D}_{\text{root}},\, G \concat \langle n \rangle_8,\, \sigaggmore,\, \sigagglast,\, \tauroot}$
    \State \Return $C_0 \concat C_1 \concat \cdots \concat C_n \concat T$
  \end{algorithmic}
\end{algorithm}

The leaf loop is annotated \textbf{in parallel} to make explicit that
the $n$ leaf computations are mutually independent; \Cref{sec:impl}
discusses the SIMD realizations. The annotation is an implementation
freedom, not a mode parameter: evaluating leaves serially, in SIMD
batches, or across cores leaves the transcript and wire format unchanged.

\paragraph{Streaming aggregation.}
The string $G = T_1 \concat \cdots \concat T_n$ in
\Cref{alg:enc,alg:dec} names the leaf tags only for specification. The
closing \Call{Absorb}{} consumes $G \concat \langle n \rangle_8$ as
$\rho$-byte duplex blocks in leaf order and advances the root state
after each, so an implementation need not hold $G$ in full. It can absorb
leaf tags in order as they become available, or buffer completed
out-of-order leaves until the next in-order tag or SIMD batch can be
absorbed, keeping only a bounded number of tags at a time. With
the bounded $1600$-bit state of every duplex, this keeps the working
memory constant in $\|M\|$ for a fixed degree of parallelism, growing
only with the number of leaves evaluated at once rather than with the
message length; \Cref{sec:impl} describes the batched realization.

\subsection{Decryption and Correctness}
\label{sec:dec}

\Cref{alg:dec} specifies $\Dec_K(U, A, C \concat T)$. Decryption mirrors
encryption at the duplex block level: each ciphertext block is absorbed under
the same phase suffix, and the plaintext block is recovered as
$m_i = c_i \oplus Z^{(i)}$ from the keystream value that encryption would have
used at that position. The edge cases are those of \Cref{sec:enc}. The
procedure absorbs all in-domain ciphertext bytes before comparing the final tag,
so the verification transcript is determined by $(K, U, A, C)$ whether or not the
submitted tag accepts.

\begin{algorithm}[t]
  \caption{$\Dec_K(U, A, C \concat T)$.}
  \label{alg:dec}
  \begin{algorithmic}[1]
    \Require $K, U, A$ as in \Cref{alg:enc}, ciphertext $C$ with $\|C\| \leq M_{\max}$, candidate tag $T \in \bits^{\tauroot}$.
    \Ensure $M$ with $\|M\| = \|C\|$, or $\bot$.
    \State $C_0 \gets$ first $\min(\|C\|, \beta)$ bytes of $C$ (possibly $\varepsilon$).
    \State Partition the remaining suffix into chunks $C_1, \dots, C_n$ of exactly $\beta$ bytes each, except that the last chunk may be shorter but nonempty, so $n = \leaves(C)$.
    \State $\mathcal{D}_{\text{root}} \gets$ new $\Duplex$
    \If{$A = \varepsilon$}
      \State $Z \gets \mathcal{D}_{\text{root}}.\duplexing(\prfxroot \concat K \concat U \concat \siginit, 8\min(\|C_0\|, \rho))$
    \Else
      \State $\mathcal{D}_{\text{root}}.\duplexing(\prfxroot \concat K \concat U \concat \siginit, 0)$
      \State $Z \gets \Call{Absorb}{\mathcal{D}_{\text{root}}, A, \sigadmore, \sigadlast, 8\min(\|C_0\|, \rho)}$
    \EndIf
    \If{$n = 0$}
      \State $(M_0,\, T') \gets \Call{StreamDec}{\mathcal{D}_{\text{root}}, C_0, Z, \sigmsgmore, \sigmsglast, \tauroot}$
      \State \Return $M_0$ if $T' = T$ (compared in constant time), else $\bot$.
    \EndIf
    \State $(M_0,\, \_) \gets \Call{StreamDec}{\mathcal{D}_{\text{root}}, C_0, Z, \sigmsgmore, \sigmsglast, 0}$
    \For{$j \gets 1$ \textbf{to} $n$ \textbf{in parallel}}
      \State $\mathcal{D}_j \gets$ new $\Duplex$
      \State $Z_j \gets \mathcal{D}_j.\duplexing(\prfxleaf \concat K \concat U \concat \langle j \rangle_8 \concat \siginit, 8\min(\|C_j\|, \rho))$
      \State $(M_j, T_j) \gets \Call{StreamDec}{\mathcal{D}_j, C_j, Z_j, \sigmsgmore, \sigmsglast, \tauleaf}$
    \EndFor
    \State $G \gets T_1 \concat \cdots \concat T_n$.
    \State $T' \gets \Call{Absorb}{\mathcal{D}_{\text{root}},\, G \concat \langle n \rangle_8,\, \sigaggmore,\, \sigagglast,\, \tauroot}$
    \If{$T' = T$ (compared in constant time)}
      \State \Return $M_0 \concat M_1 \concat \cdots \concat M_n$
    \Else
      \State \Return $\bot$
    \EndIf
  \end{algorithmic}
\end{algorithm}

\paragraph{Correctness.}
For every $K \in \bits^k$, $U \in \bits^{\nu}$,
$A$ with $\|A\| \leq A_{\max}$, and $M$ with $\|M\| \leq M_{\max}$,
$\Dec_K(U, A, \Enc_K(U, A, M)) = M$. Encryption and decryption absorb the same
ciphertext blocks under the same duplex instances and phase suffixes, so their
duplex states and final root tags agree. Each recovered block satisfies
$c_i \oplus Z^{(i)} = (m_i \oplus Z^{(i)}) \oplus Z^{(i)} = m_i$.

\begin{lemma}[\TW{} tidiness]
\label{lem:tw-tidiness}
For every valid \TW{} decryption input, if
$\Dec_K(U, A, C \concat T) = M \neq \bot$, then
$\Enc_K(U, A, M) = C \concat T$.
\end{lemma}

\begin{proof}
An accepting decryption has $\|M\| = \|C\|$ and therefore the same canonical
chunking as re-encryption of $M$. The decryption transcript absorbs the same
$(K, U, A)$ values and ciphertext blocks that this re-encryption absorbs, under
the same duplex instances and phase suffixes. Its stream outputs recover
$m_i = c_i \oplus Z^{(i)}$, so re-encryption emits the submitted ciphertext blocks
again. The recomputed leaf tags and final root candidate tag are therefore the
same values; since decryption accepts only when the final candidate equals the
submitted tag $T$, re-encryption returns $C \concat T$.
\end{proof}
